% ============================================================================
% Fichier  : partie2/P02_C10_collecter_preuve_nek.tex
% Titre    : Collecter la Preuve : l'Instrument OUI/NON/NEK
% Auteur   : MINKA MI NGUIDJOÏ Thierry Emmanuel
% Version  : 1.0 -- novembre 2026
% ============================================================================
% Base : ACSI C08 (Ω ; ρ = 0,80 - 0,65 = 0,15 ; hiérarchie de fiabilité ;
% types DOC, ENT, OBS, TEST, MIX ; n ≈ F/t ; D.2.1 31 départs 19/7/5 ;
% lettre d'extraction GUID 5100 ; MEFOU).
% Sources : A0, A1 (hal-05701270) ; INTOSAI GUID 5100 ; échantillonnage par
% attributs : loi binomiale et approximation de Poisson (facteurs de
% fiabilité 2,31 / 3,89 à 90 %, 3,00 / 4,74 à 95 %) ; échantillonnage en
% unités monétaires (méthode classique). Exemple chiffré sur les 4 812
% paiements de 2025 du lab ACSI-LAB-07 : total 6 363 640 000 FCFA, plus
% gros paiement 22 784 000 FCFA (calcul du 06/11/2026).
% Corrigés : corriges/Corr_C10.tex.
% ============================================================================

\chapter[Collecter la Preuve]{Collecter la Preuve : l'Instrument OUI/NON/NEK}
\label{chap:preuve}
\chaptermark{Collecter la preuve}

\epigraph{%
    « Ce qui n'est pas dans le dossier\\
    n'a pas été audité. »%
}{Règle des papiers de travail}

\niveauChapitre{Intermédiaire}{%
    Maîtrise du Chapitre~\ref{chap:decomposition} (gabarit canonique, fiche
    d'assertion) et du Chapitre~\ref{chap:risque} (risque de détection)%
}

\begin{boxobjectifs}
\begin{itemize}
    \item Justifier l'instrument de réponse \OUI{}, \NON{} ou \NEK{} et le
          distinguer de la réponse « non applicable ».
    \item Mesurer le biais du choix forcé d'une grille binaire.
    \item Classer les preuves selon leur source et leur fiabilité.
    \item Conduire un entretien d'audit à partir des fiches d'assertion.
    \item Démontrer la formule de taille d'échantillon par attributs,
          l'appliquer avec ou sans écart attendu, et conclure sur un
          échantillon.
    \item Dimensionner un échantillonnage en unités monétaires.
    \item Garantir l'intégrité des extractions et tenir des papiers de
          travail jusqu'à la trace.
\end{itemize}
\end{boxobjectifs}

\bigskip

% ============================================================================
\section*{Question directrice : comment savoir qu'on sait~?}
% ============================================================================

La décomposition du domaine~D est prête (Chapitre~\ref{chap:decomposition}).
Pour chaque contrôle, une question attend sa réponse. M.~BEYALA pose à son
équipe une règle simple~: une réponse ne vaut que par la preuve qui
l'appuie, et l'on a le droit, et même le devoir, de répondre « je ne sais
pas » quand la preuve manque.

\medskip
\textbf{Plan du chapitre.}
\begin{enumerate}
    \item L'instrument \OUI{}, \NON{}, \NEK{} et le biais du choix forcé
          (sections~\ref{sec:c10-instrument} et~\ref{sec:c10-biais}).
    \item Les preuves et leur fiabilité (section~\ref{sec:c10-preuves}).
    \item Les entretiens (section~\ref{sec:c10-entretiens}).
    \item Les tests (section~\ref{sec:c10-tests}).
    \item L'échantillonnage par attributs, démontré
          (section~\ref{sec:c10-attributs}), et en unités monétaires
          (section~\ref{sec:c10-mus}).
    \item Les extractions et les papiers de travail
          (section~\ref{sec:c10-papiers}).
\end{enumerate}

\separateur

% ============================================================================
\section{L'Instrument OUI / NON / NEK}
\label{sec:c10-instrument}
% ============================================================================

\subsection{Trois réponses}

L'auditologie établit, à partir de ses axiomes d'observabilité partielle, de
cognition limitée, de mesurabilité et d'interaction stratégique, que tout
instrument de réponse d'audit valide doit comporter une modalité explicite
de connaissance insuffisante~\cite{minka2026auditologie}. L'article A1 en
construit l'instrument et en étudie les propriétés~\cite{minka2026nek}.

\begin{boxdef}
\termecle{Espace de réponse trichotomique}\index{espace de réponse
trichotomique}~: ensemble $\Omega = \{\OUI, \NON, \NEK\}$ des réponses
admises à une question d'audit.
\begin{itemize}
    \item \OUI{}~: le critère est satisfait, et une preuve l'établit.
    \item \NON{}~: le critère n'est pas satisfait, et une preuve l'établit.
    \item \NEK{} (\textit{Not Enough Knowledge})~: l'information disponible
          ne permet pas de conclure, malgré des diligences normales.
\end{itemize}
\end{boxdef}

\subsection{Deux usages}

L'instrument sert à deux moments de la mission.
\begin{enumerate}
    \item \textbf{Auprès des répondants}~: une même question est posée à
          plusieurs interlocuteurs (entretiens, questionnaires). Pour chaque
          question~$q$, on obtient une distribution de réponses, dont l'article
          A1 tire trois mesures~: $\mu_q$, la part des \OUI{}~; $U_q$, la part
          des \NEK{}~; $H_q$, l'entropie~\cite{minka2026nek}.
    \item \textbf{Par l'auditeur}~: après ses tests, l'auditeur cote chaque
          contrôle \OUI{}, \NON{} ou \NEK{}, preuve à l'appui. C'est la cote
          qui entre dans le rapport.
\end{enumerate}
Les deux usages se répondent~: une question où les répondants divergent
(entropie élevée) appelle des tests~; les tests permettent ensuite à
l'auditeur de coter avec preuve.

\subsection{NEK n'est pas « non applicable »}

\begin{boxalert}
\textbf{Violation typique.} Coter \NEK{} un contrôle qui ne s'applique pas au
périmètre (par exemple, la double authentification d'un service qui n'est
pas exposé à Internet).

\smallskip
\textbf{Règle~:} un contrôle sans objet dans le périmètre est \emph{non
applicable}~: il sort de la décomposition, avec la raison écrite. \NEK{} est
réservé à un contrôle applicable dont on ne peut pas établir l'état. Les
confondre fausse le taux d'incertitude~$U$.
\end{boxalert}

% ============================================================================
\section{Le Biais du Choix Forcé}
\label{sec:c10-biais}
% ============================================================================

Une grille binaire force le répondant qui ne sait pas à choisir. Il choisit
le plus souvent la réponse attendue, ou la plus confortable~: le \OUI{}. La
grille binaire \emph{surestime} donc la conformité. L'article A1 mesure cette
surestimation par un indice de réduction du biais~\cite{minka2026nek}~:
\[
    \rho_q = \hat{\mu}_q^{B} - \mu_q^{T},
\]
où $\hat{\mu}_q^{B}$ est la part des \OUI{} obtenue avec l'instrument binaire et
$\mu_q^{T}$ celle obtenue avec l'instrument trichotomique, sur les mêmes
questions.

\begin{boxlecture}{Domaine D, binaire puis trichotomique}
Lors de la première mission, les quatre interlocuteurs avaient répondu aux
cinq questions du domaine~D avec une grille binaire~: 16~\OUI{} sur 20
réponses. Au second passage, avec \NEK{}~: 13~\OUI{}, 2~\NON{}, 5~\NEK{}.

\medskip
\textbf{Constatation.} $\hat{\mu}^{B} = 16/20 = 0{,}80$~; $\mu^{T} = 13/20 =
0{,}65$~; $\rho = 0{,}15$.

\medskip
\textbf{Analyse} (confiance moyenne). La grille binaire surestimait la
conformité du domaine de quinze points~: trois répondants qui ne savaient
pas avaient répondu \OUI{}. C'est ainsi que le premier rapport a pu conclure
à un domaine «~globalement maîtrisé~».

\medskip
\hfill\lecture{0{,}65}{0{,}25}{1{,}24}
\end{boxlecture}

% ============================================================================
\section{Les Preuves et leur Fiabilité}
\label{sec:c10-preuves}
% ============================================================================

\subsection{Suffisante et appropriée}

\begin{boxdef}
\termecle{Preuve d'audit}\index{preuve d'audit}~: information sur laquelle
l'auditeur fonde sa cotation et ses conclusions. Elle doit être
\emph{suffisante} (en quantité, au regard du risque) et \emph{appropriée}
(pertinente pour le critère, et fiable).
\end{boxdef}

\subsection{Une hiérarchie de fiabilité}

\begin{tableaucours}{cL{4.2cm}L{4.4cm}L{2.8cm}}{%
    \tableauHeader{Rang} &
    \tableauHeader{Source} &
    \tableauHeader{Exemple} &
    \tableauHeader{Fiabilité}
}
1 & Trace produite par le système, extraite par l'auditeur & Journal des
désactivations extrait en sa présence & La plus élevée \\
2 & Document externe & Relevé bancaire, attestation d'un prestataire &
Élevée \\
3 & Document interne approuvé & Compte rendu de revue signé & Moyenne \\
4 & Observation directe & L'auditeur voit le magasinier saisir une
réception & Ponctuelle \\
5 & Déclaration orale & « Les comptes sont revus chaque trimestre » & La
plus faible \\
\end{tableaucours}
\vspace{-0.8em}
\begin{center}{\small\itshape Tableau~10.1. Fiabilité des preuves selon leur
source}\end{center}

Une déclaration orale ne suffit jamais à coter \OUI{}~: elle oriente la
recherche de la preuve. Lorsque l'auditeur ne dispose que d'une
déclaration, la cote est \NEK{}, et la question devient «~où est la
preuve~?~».

\subsection{La corroboration}

Deux preuves indépendantes qui concordent valent plus que la plus fiable
des deux prise seule. L'auditeur cherche donc, pour les contrôles
importants, une seconde source~: une revue signée (rang~3) recoupée avec les
dates de modification des droits dans le journal (rang~1). Si les deux
sources se contredisent sans qu'on puisse les départager, la cote est
\NEK{}, et l'écart lui-même devient un objet d'investigation.

% ============================================================================
\section{Les Entretiens}
\label{sec:c10-entretiens}
% ============================================================================

Le guide d'entretien (outil~O5) dérive des fiches d'assertion~: chaque
question de fiche devient une question d'entretien, classée par
interlocuteur. L'entretien d'audit, à la différence de l'entretien de prise
de connaissance (Chapitre~\ref{chap:planifier}), recueille des réponses
cotables~; il obéit à des règles.

\begin{boxmetier}{Auditeur SI}
\textbf{Conduire un entretien d'audit}
\begin{enumerate}
    \item Annoncer l'objet, la durée, et que la réponse \NEK{} est admise.
    \item Commencer par des questions ouvertes («~Comment se passe le départ
          d'un agent~?~»), puis passer aux questions de la fiche.
    \item Ne pas suggérer la réponse attendue~: «~Les comptes sont-ils
          désactivés sous 48~heures~?~» plutôt que «~Vous désactivez bien
          les comptes, n'est-ce pas~?~».
    \item Demander, pour chaque réponse affirmative, la pièce qui l'établit.
    \item Mener l'entretien à deux, l'un interrogeant, l'autre notant.
    \item Rédiger le compte rendu le jour même et le faire valider par
          l'interlocuteur.
\end{enumerate}
\end{boxmetier}

\begin{boiteRemarque}{Le répondant peut-il se cacher derrière NEK~?}
Oui~: le Chapitre~\ref{chap:fraude} l'a montré, le taux de \NEK{} peut
augmenter avec l'enjeu. C'est pourquoi l'auditeur note le contexte de
chaque réponse (entretien formel, questionnaire, auto-évaluation) et
compare les contextes. Un \NEK{} n'est jamais une fin~: il appelle un test
qui cherche la preuve directement dans le système.
\end{boiteRemarque}

% ============================================================================
\section{Tester}
\label{sec:c10-tests}
% ============================================================================

\begin{tableaucours}{cL{3.0cm}L{4.4cm}L{3.6cm}}{%
    \tableauHeader{Code} &
    \tableauHeader{Type} &
    \tableauHeader{Principe} &
    \tableauHeader{Exemple (domaine D)}
}
DOC & Revue documentaire & Lire une procédure, un compte rendu, un contrat
& Lire la procédure de départ \\
ENT & Entretien & Recueillir une déclaration, à confirmer & Interroger le
responsable RH \\
OBS & Observation & Voir le contrôle s'exécuter & Assister à une revue des
comptes \\
TEST & Test technique & Ré-exécuter le contrôle ou analyser les données &
Rapprocher départs et comptes actifs \\
MIX & Combinaison & Plusieurs types sur un même contrôle & Revue, puis
test \\
\end{tableaucours}
\vspace{-0.8em}
\begin{center}{\small\itshape Tableau~10.2. Les types de tests du plan de
tests (outil O7)}\end{center}

Un test peut viser la \emph{conception} d'un contrôle (test de cheminement
sur un cas), son \emph{efficacité} (échantillon sur la période) ou
directement les \emph{données} (tests substantifs,
Chapitre~\ref{chap:audit-applications}).

\begin{boxscenario}{Opération LIZBIZ. Tester D.2.1}
\textbf{Contrôle.} Désactivation des comptes des agents partis sous
48~heures (fiche du prologue).

\medskip
\textbf{Choix.} Les départs sont peu nombreux (31 sur la période)~: AUDIX
CONSEIL teste la population entière plutôt qu'un échantillon, par un
rapprochement entre la liste des départs des ressources humaines et le
journal des désactivations.

\medskip
\textbf{Constatation.} 31~départs~; 19 comptes désactivés sous 48~heures,
7 désactivés plus tard, 5 encore actifs.

\medskip
\textbf{Cotation.} \NON{}~: le critère n'est pas satisfait, et la preuve est
établie sur la population entière.
\end{boxscenario}

\begin{boxoutils}{Le rapprochement de D.2.1}
\begin{lstlisting}[language=SQL]
-- departs(matricule, date_depart) : ressources humaines
-- comptes(login, matricule, date_desactivation) : application
SELECT d.matricule, d.date_depart, c.login, c.date_desactivation,
  CASE
    WHEN c.date_desactivation IS NULL THEN 'ACTIF'
    WHEN julianday(c.date_desactivation)
         - julianday(d.date_depart) <= 2 THEN 'SOUS 48 H'
    ELSE 'EN RETARD'
  END AS statut
FROM departs d LEFT JOIN comptes c ON c.matricule = d.matricule;
\end{lstlisting}
Sur la période, la requête classe les 31~départs en 19 « sous 48~h »,
7 « en retard » et 5 « actif ». La jointure externe (\texttt{LEFT JOIN})
fait aussi apparaître un départ sans compte, à examiner à part.
\end{boxoutils}

% ============================================================================
\section{L'Échantillonnage par Attributs}
\label{sec:c10-attributs}
% ============================================================================

\subsection{Trois méthodes}

On ne teste pas toutes les exécutions d'un contrôle~: on en tire un
\termecle{échantillon}\index{échantillon}. Trois méthodes coexistent.
\begin{itemize}
    \item \textbf{Par jugement}~: l'auditeur choisit les cas les plus risqués
          (les plus gros montants, les dernières semaines de la période). Il
          ne peut pas extrapoler, mais il cible.
    \item \textbf{Aléatoire}~: chaque élément a la même chance d'être tiré~;
          l'auditeur conserve la méthode et la graine du tirage pour pouvoir
          le refaire.
    \item \textbf{Statistique par attributs}~: on fixe un niveau de confiance
          et un taux d'écart tolérable, et l'on en déduit la taille de
          l'échantillon.
\end{itemize}

\subsection{La formule, démontrée}

Soit une population où une proportion $p$ des exécutions du contrôle sont
en écart (le contrôle n'a pas fonctionné). On tire au hasard $n$
exécutions. Si la population est grande, le nombre d'écarts dans
l'échantillon suit une loi binomiale de paramètres $n$ et $p$~; la
probabilité de ne trouver \emph{aucun} écart est
\[
    P(0) = (1 - p)^n.
\]
L'auditeur veut pouvoir conclure, s'il ne trouve aucun écart, que $p$ ne
dépasse pas un taux tolérable $t$, avec un risque $\beta$ de se tromper
(par exemple $\beta = 5~\%$ pour une confiance de 95~\%). Il faut donc que,
si le vrai taux valait $t$, la probabilité de ne rien trouver soit au plus
$\beta$~:
\[
    (1 - t)^n \le \beta
    \quad\Longleftrightarrow\quad
    n \ge \frac{\ln \beta}{\ln(1 - t)}.
\]
Pour $t$ petit, $\ln(1 - t) \approx -t$, d'où l'approximation
\[
    n \approx \frac{-\ln \beta}{t} = \frac{F}{t},
    \qquad F = \ln\frac{1}{\beta},
\]
avec $F = \ln 10 \approx 2{,}30$ à 90~\% et $F = \ln 20 \approx 3{,}00$ à
95~\%. C'est l'approximation de Poisson~: le nombre d'écarts d'un grand
échantillon où $p$ est petit suit approximativement une loi de Poisson de
moyenne $np$.

\begin{boxinfo}
\textbf{Exact ou approché~?} Pour $t = 5~\%$ et 95~\% de confiance, la
formule exacte donne $n \ge \ln 0{,}05 / \ln 0{,}95 \approx 58{,}4$, soit
59~exécutions~; l'approximation donne $3{,}00 / 0{,}05 = 60$. Pour $t =
1~\%$ et un risque de 10~\%, $n \ge \ln 0{,}1 / \ln 0{,}99 \approx 229{,}1$,
soit les 230~éléments évoqués au Chapitre~\ref{chap:audit-auditologie}.
L'approximation est légèrement prudente~; elle suffit en pratique.
\end{boxinfo}

\subsection{Si l'on accepte un écart}

Lorsque l'auditeur accepte de trouver un ou deux écarts sans conclure à
l'inefficacité, le facteur $F$ augmente. Il se lit dans la table de la loi
de Poisson~: c'est la plus grande moyenne $\lambda$ telle que la
probabilité d'observer au plus $k$ écarts soit encore égale à $\beta$.

\begin{tableaucours}{L{4.2cm}ccc}{%
    \tableauHeader{Écarts acceptés $k$} &
    \tableauHeader{0} & \tableauHeader{1} & \tableauHeader{2}
}
$F$ à 90~\% de confiance & 2,31 & 3,89 & 5,33 \\
$F$ à 95~\% de confiance & 3,00 & 4,74 & 6,30 \\
\end{tableaucours}
\vspace{-0.8em}
\begin{center}{\small\itshape Tableau~10.3. Facteurs de fiabilité (loi de
Poisson)}\end{center}

\begin{tableaucours}{L{4.0cm}cccc}{%
    \tableauHeader{Taux tolérable $t$} &
    \tableauHeader{90~\%, $k=0$} & \tableauHeader{95~\%, $k=0$} &
    \tableauHeader{90~\%, $k=1$} & \tableauHeader{95~\%, $k=1$}
}
10~\% & 23 & 30 & 39 & 47 \\
5~\% & 46 & 60 & 78 & 95 \\
2~\% & 116 & 150 & 195 & 237 \\
\end{tableaucours}
\vspace{-0.8em}
\begin{center}{\small\itshape Tableau~10.4. Tailles d'échantillon $n =
F/t$, arrondies à l'entier le plus proche}\end{center}

\subsection{Conclure sur l'échantillon}

Après le test, l'auditeur calcule la \textbf{borne supérieure} du taux
d'écart~: $p_{\max} = F_k / n$, où $k$ est le nombre d'écarts trouvés. Il la
compare au taux tolérable~: si $p_{\max} \le t$, le contrôle est efficace
au niveau de confiance choisi~; sinon, il ne l'est pas.

\begin{boxscenario}{Opération LIZBIZ. Les validations de paiements}
\textbf{Test.} Le contrôle « validation par le DAF de tout paiement » est
testé sur 60~paiements tirés au hasard (95~\%, $t = 5~\%$, aucun écart
attendu). Un paiement de l'échantillon n'a pas de validation.

\medskip
\textbf{Calcul.} $p_{\max} = F_1 / n = 4{,}74 / 60 \approx 7{,}9~\%$, plus
que les 5~\% tolérables.

\medskip
\textbf{Cotation.} \NON{}~: on ne peut pas conclure, avec 95~\% de
confiance, que le taux d'écart est inférieur à 5~\%. Pour conclure à
l'efficacité malgré cet écart, il aurait fallu un échantillon de 95
(Tableau~10.4) sans nouvel écart.
\end{boxscenario}

\begin{boxalert}
\textbf{Violation typique.} Trouver un écart, le qualifier « d'isolé » et
conclure que le contrôle est efficace.

\smallskip
\textbf{Règle~:} la taille d'échantillon a été calculée pour zéro écart~;
un écart invalide la conclusion prévue. L'auditeur élargit l'échantillon
selon le Tableau~10.4, ou conclut à l'inefficacité. Il recherche aussi la
cause de l'écart~: un écart intentionnel n'est jamais isolé.
\end{boxalert}

Lorsque les données sont extractibles, il vaut souvent mieux tester
l'ensemble de la population, ce qui supprime le risque d'échantillonnage
(Chapitre~\ref{chap:audit-continu}).

% ============================================================================
\section{L'Échantillonnage en Unités Monétaires}
\label{sec:c10-mus}
% ============================================================================

\subsection{Le principe}

Pour tester l'exactitude de montants (une population de paiements, de
factures), on tire non pas des opérations, mais des \emph{francs}~: chaque
franc a la même chance d'être tiré, et l'opération qui le contient est
examinée. Les grosses opérations ont donc plus de chances d'être retenues,
ce qui concentre l'effort là où l'enjeu est le plus élevé.

\begin{boxdef}
\termecle{Échantillonnage en unités monétaires}\index{échantillonnage en
unités monétaires}~: méthode qui tire un échantillon d'unités monétaires
dans une population de valeur comptable $V$, avec
\[
    n = \frac{V \times F}{S}, \qquad J = \frac{V}{n},
\]
où $S$ est l'anomalie tolérable (seuil de signification retenu pour le
test), $F$ le facteur de fiabilité (Tableau~10.3) et $J$ l'intervalle de
tirage. Toute opération de montant supérieur à $J$ est retenue d'office.
\end{boxdef}

\begin{boxscenario}{Opération LIZBIZ. Les paiements de 2025}
\textbf{Données.} 4~812~paiements, pour une valeur totale de
6~363~640~000~FCFA~; le plus gros paiement vaut 22~784~000~FCFA.

\medskip
\textbf{Paramètres.} Anomalie tolérable $S = 63{,}6$~millions (1~\% de la
valeur), confiance 95~\%, aucune anomalie attendue ($F = 3{,}00$).

\medskip
\textbf{Calcul.}
\[
n = \frac{6~363{,}64 \times 3{,}00}{63{,}6} \approx 300,
\qquad J = \frac{6~363{,}64}{300} \approx 21{,}2~\text{millions FCFA}.
\]
Les paiements supérieurs à 21,2~millions (dont celui de 22,8~millions)
sont retenus d'office~; les autres sont tirés en parcourant la population
cumulée par pas de 21,2~millions, à partir d'un point de départ aléatoire.

\medskip
\textbf{Ce que la méthode ne voit pas.} Les 23~paiements de 4,9~millions
de TRANSLOG SERVICES ont chacun une probabilité d'environ $4{,}9/21{,}2
\approx 23~\%$ d'être tirés~: on en attend environ cinq dans
l'échantillon. Mais ces paiements sont \emph{exacts}~: leur montant
correspond à la facture. L'échantillonnage monétaire teste l'exactitude,
pas la réalité de la prestation~; la fraude du fournisseur fictif se
détecte par les requêtes du Chapitre~\ref{chap:fraude}.
\end{boxscenario}

% ============================================================================
\section{Extractions et Papiers de Travail}
\label{sec:c10-papiers}
% ============================================================================

\subsection{Garantir l'intégrité des extractions}

Une extraction de données remise par l'audité n'est une preuve que si l'on
peut établir d'où elle vient et qu'elle n'a pas été modifiée. Le GUID~5100
de l'INTOSAI recommande qu'elle soit accompagnée d'une lettre qui en
précise la source, l'horodatage, l'empreinte numérique, les paramètres
d'extraction et l'environnement d'origine~\cite{intosai2019guid5100}.

\begin{boxoutils}{Calculer et vérifier une empreinte}
\textbf{Calculer l'empreinte d'une extraction} (Linux, ou terminal de
NDAP SUKULU)~:

\texttt{sha256sum comptes\_actifs\_2026-04-30.csv > comptes.sha256}

\medskip
\textbf{Vérifier plus tard qu'elle n'a pas changé}~:

\texttt{sha256sum -c comptes.sha256}

\medskip
Résultat attendu~: \texttt{comptes\_actifs\_2026-04-30.csv: OK}. Toute
modification du fichier, même d'un caractère, change l'empreinte. Sous
Windows, la commande \texttt{certutil -hashfile fichier SHA256} donne le
même résultat.
\end{boxoutils}

\begin{boxscenario}{Opération MEFOU. La lettre d'extraction}
\textbf{Faits.} La Chambre des comptes reçoit de l'ANPP l'extraction des
comptes du système de paiement.

\medskip
\textbf{Contenu de la lettre} signée par la cheffe de la division
informatique~: base de production, requête utilisée, date et heure
d'extraction, nombre de lignes, empreinte SHA-256 du fichier.

\medskip
\textbf{Analyse} (confiance élevée). Le magistrat peut prouver que les écarts
qu'il relèvera proviennent bien des données de l'agence, et non d'un fichier
retouché~: la constatation devient opposable.
\end{boxscenario}

\begin{boxalert}
\textbf{Violation typique.} Ouvrir l'extraction dans un tableur, la trier,
l'enregistrer sous le même nom, puis calculer l'empreinte.

\smallskip
\textbf{Règle~:} l'empreinte se calcule sur le fichier reçu, avant toute
ouverture~; l'auditeur travaille sur une copie. L'original et son empreinte
restent intacts dans le dossier.
\end{boxalert}

\subsection{Le papier de travail d'un test}

\begin{boxprocedure}
\textbf{Le papier de travail d'un test.} Pour qu'un tiers puisse refaire le
test et aboutir à la même cotation, chaque papier de travail indique~:
\begin{itemize}
    \item sa référence, et celle du contrôle dans la décomposition (par
          exemple PT-14, contrôle D.2.1)~;
    \item l'objectif du test, le critère et la population ou l'échantillon,
          avec sa méthode de sélection (et la graine, si le tirage est
          aléatoire)~;
    \item les preuves examinées, avec leur source et leur empreinte, et les
          traces correspondantes~;
    \item les résultats, les écarts relevés, la borne supérieure du taux
          d'écart s'il y a échantillonnage, et la cotation justifiée~;
    \item le nom de l'auditeur, la date, et le visa du réviseur.
\end{itemize}
\end{boxprocedure}

\begin{boxinfo}
\textbf{Le commissaire aux comptes.} Le commissaire aux comptes de LiZbiZ
peut utiliser les travaux d'AUDIX CONSEIL sur le contrôle d'accès pour
fonder sa propre opinion sur les comptes. Il ne le fera que si les papiers
de travail lui permettent d'apprécier la qualité des preuves~: méthode
d'échantillonnage, intégrité des extractions, cotations justifiées. Un
\NEK{} documenté lui est plus utile qu'un \OUI{} sans preuve.
\end{boxinfo}

\begin{boxmetier}{Auditeur de sécurité}
\textbf{Extraire soi-même, ou faire extraire devant soi.} Pour les
journaux de sécurité, l'auditeur de sécurité fournit la commande ou la
requête, assiste à l'extraction, calcule l'empreinte sur place et la fait
contresigner par l'administrateur. Une extraction remise plus tard par
courriel n'a que la valeur d'un document interne.
\end{boxmetier}

\begin{boxinfo}
\textbf{Ce que dit la recherche.} L'article A1~\cite{minka2026nek} montre
que l'ajout de \NEK{} ne fait pas qu'ajouter une case~: il change ce que
mesure l'instrument, en séparant l'ignorance ($U$) de la non-conformité.
L'échantillonnage statistique fait la même chose pour l'auditeur~: la borne
supérieure $p_{\max}$ dit ce que l'échantillon permet d'affirmer, et rien de
plus.
\end{boxinfo}

\separateur

\begin{boxresume}
\textbf{Ce qu'il faut retenir de ce chapitre~:}
\begin{itemize}
    \item \textbf{Trois réponses}~: \OUI{} et \NON{} s'appuient sur une
          preuve~; \NEK{} déclare une connaissance insuffisante~; « non
          applicable » sort du périmètre et n'est pas \NEK{}.
    \item \textbf{Le choix forcé surestime la conformité}~: l'indice
          $\rho = \hat{\mu}^{B} - \mu^{T}$ mesure ce biais.
    \item \textbf{Les preuves n'ont pas la même fiabilité}~: la trace extraite
          par l'auditeur l'emporte sur la déclaration orale, qui ne suffit
          jamais à coter \OUI{}~; deux preuves concordantes valent plus
          qu'une.
    \item \textbf{L'entretien} suit les fiches d'assertion, sans suggérer la
          réponse, et demande la pièce.
    \item \textbf{L'échantillonnage par attributs} repose sur $(1-t)^n \le
          \beta$, d'où $n \approx F/t$~; un écart trouvé change $F$ et peut
          invalider la conclusion.
    \item \textbf{L'échantillonnage monétaire} concentre l'effort sur les
          gros montants et teste l'exactitude, pas la réalité.
    \item \textbf{L'intégrité des extractions} se prouve par leur empreinte~;
          les papiers de travail permettent à un tiers de refaire chaque test.
\end{itemize}
\end{boxresume}

\begin{boxinfo}
\textbf{Points de contrôle de l'élève.} Je sais~: distinguer \NEK{} et non
applicable~; calculer $\rho$~; classer une preuve~; démontrer $n \approx
F/t$~; calculer une borne $p_{\max}$~; dimensionner un échantillon
monétaire~; calculer une empreinte SHA-256.
\end{boxinfo}

\bigskip

\begin{boxexo}[title={\ding{45}\quad Exercice~10.1. OUI, NON, NEK ou non applicable~? \niveaubase}]
Cotez chaque situation et justifiez.
\begin{enumerate}[(a)]
    \item Le responsable informatique affirme que les journaux sont revus
          chaque semaine, sans pouvoir le montrer.
    \item Le journal des désactivations montre que tous les départs du
          trimestre ont été traités sous 48~heures.
    \item Le contrôle de la double authentification à distance, alors
          qu'aucun accès à distance n'existe.
    \item Le journal de revue des anomalies existe, mais il est vide pour
          les trois derniers mois, alors que la procédure exige une revue
          hebdomadaire.
\end{enumerate}
\end{boxexo}

\begin{boxexo}[title={\ding{45}\quad Exercice~10.2. Le biais \niveaubase}]
Sur une question posée à dix répondants, la grille binaire donnait
8~\OUI{}~; la grille trichotomique donne 5~\OUI{}, 1~\NON{}, 4~\NEK{}.
Calculez $\rho$, $U$ et $H$, et interprétez.
\end{boxexo}

\begin{boxexo}[title={\ding{45}\quad Exercice~10.3. Classer des preuves \niveaubase}]
Classez par fiabilité décroissante, pour le contrôle D.4.2~: (a) le
compte rendu signé de la revue du premier trimestre~; (b) l'affirmation du
RSSI~; (c) l'extraction des comptes d'administration faite devant
l'auditeur~; (d) un courriel du DSI disant «~revue faite~». Quelle
combinaison vous permettrait de coter \OUI{}~?
\end{boxexo}

\begin{boxexo}[title={\ding{45}\quad Exercice~10.4. Démontrer \niveaumoyen}]
\begin{enumerate}[(a)]
    \item Calculez la taille exacte et la taille approchée de l'échantillon
          pour $t = 10~\%$ et 90~\% de confiance.
    \item Montrez que $\ln(1-t) \approx -t$ en utilisant le développement
          $\ln(1-t) = -t - t^2/2 - \dots$ et expliquez pourquoi
          l'approximation est prudente.
\end{enumerate}
\end{boxexo}

\begin{boxexo}[title={\ding{45}\quad Exercice~10.5. Dimensionner et conclure \niveaumoyen}]
\begin{enumerate}[(a)]
    \item Quelle taille d'échantillon retenir pour tester, avec une confiance
          de 95~\% et un taux d'écart tolérable de 5~\%, la validation des
          ajustements de stock de LiZbiZ~?
    \item Le test ne révèle aucun écart. Calculez la borne supérieure du
          taux d'écart et concluez.
    \item Le test révèle deux écarts. Calculez la borne et concluez. Que
          doit faire l'auditeur~?
\end{enumerate}
\end{boxexo}

\begin{boxexo}[title={\ding{45}\quad Exercice~10.6. Échantillonnage monétaire \niveaumoyen}]
La paie de l'ANPP compte 12~400~versements pour 18,6~milliards de FCFA en
2025. Le magistrat retient une anomalie tolérable de 186~millions et une
confiance de 90~\%, sans anomalie attendue.
\begin{enumerate}[(a)]
    \item Calculez $n$ et l'intervalle de tirage $J$.
    \item Un versement de 410~millions figure dans la population. Que
          devient-il~?
    \item Pourquoi cette méthode ne détecterait-elle pas un salarié fantôme
          payé correctement chaque mois~?
\end{enumerate}
\end{boxexo}

\begin{boxexo}[title={\ding{45}\quad Exercice~10.7. La requête de D.2.1 \niveaumoyen}]
Modifiez la requête du chapitre pour qu'elle affiche aussi les comptes
\emph{utilisés} après la date de départ (table
\texttt{connexions(login, horodatage)}). Pourquoi cette information
change-t-elle la gravité du constat~?
\end{boxexo}

\begin{boxexo}[title={\ding{45}\quad Exercice~10.8. L'empreinte \niveaubase}]
Un auditeur reçoit \texttt{droits.csv} par courriel, l'ouvre, supprime deux
colonnes inutiles, l'enregistre, puis calcule l'empreinte. Qu'a-t-il
prouvé~? Décrivez la bonne séquence.
\end{boxexo}

\begin{boxexo}[title={\ding{45}\quad Exercice~10.9. Le plan de tests du sous-domaine D.4 \niveauexpert}]
\begin{enumerate}[(a)]
    \item À partir de la fiche d'assertion D.4.2 (Chapitre~\ref{chap:decomposition}),
          rédigez trois tests de types différents.
    \item Pour chacun, précisez la preuve attendue, sa fiabilité et la trace
          recherchée.
    \item Rédigez trois questions du guide d'entretien destinées au
          responsable de la sécurité.
    \item Rédigez la lettre d'extraction que vous demanderez pour la liste
          des comptes d'administration.
\end{enumerate}
\end{boxexo}

\begin{boxexo}[title={\ding{45}\quad Exercice~10.10. Population ou échantillon~? \niveauexpert}]
Pour chacun des contrôles suivants, dites si vous testez la population ou
un échantillon, et justifiez par le coût, le risque et la disponibilité des
données~: (a) revue trimestrielle des droits~; (b) validation de
4~812~paiements~; (c) inventaire de 3~000~références dans quatre
entrepôts~; (d) 18~340~expéditions rapprochées des factures.
\end{boxexo}

\bigskip

\begin{boxinfo}
\textbf{Pour aller plus loin.}
\begin{itemize}
    \item A1, \textit{Modelling Epistemic Uncertainty in IS
          Audit}~\cite{minka2026nek}~: l'instrument \OUI{}, \NON{}, \NEK{}
          et ses mesures.
    \item INTOSAI, GUID~5100~\cite{intosai2019guid5100}~: les preuves et
          l'intégrité des données dans l'audit des SI.
    \item IAASB, ISA~315 révisée~\cite{iaasb2019isa315}~: l'appréciation des
          risques qui fixe l'étendue des travaux.
\end{itemize}
\end{boxinfo}

\begin{boxinfo}
\textbf{Transition.} Les réponses et les preuves sont collectées. Il faut
maintenant conclure~: calculer les mesures de chaque domaine, décider ce
qu'elles permettent d'affirmer, et rédiger les observations et le rapport.
C'est l'objet du Chapitre~\ref{chap:conclure}.
\end{boxinfo}
